\everymath{\displaystyle}

%%%%%%%%%%%%%%%%%%%%%%%%%%%%%%%%%%%%%%%%
%           main body
%%%%%%%%%%%%%%%%%%%%%%%%%%%%%%%%%%%%%%%%

%-----------------------------------
%	TITLE PAGE
%-----------------------------------

\title[第三部分\\
模的定义与性质]{\texorpdfstring{\LARGE \bfseries 第三部分\quad 模的定义与性质}{第三部分 模的定义与性质}}
\author[抽象代数 2]{\Large \bfseries 抽象代数 2}
\date{}

%------------------------------------------------

\begin{document}

%------------------------------------------------

\begin{frame}

\titlepage % Print the title page as the first slide

\end{frame}

%------------------------------------------------

\section[定义]{定义: 左 R-模}

%------------------------------------------------

\begin{frame}{定义: 左 $R$-模}

\begin{Def}[左 $R$-模]
$R$ 为一个幺环, $M$ 为一个 Abel 群, 称 $M$ 为 $R$ 的一个左-模, 如果存在
$f: R \times M \to M$, 满足:
\begin{itemize}
\item[(i)] $(a_1 + a_2)\,\alpha = a_1 \alpha + a_2 \alpha$,
  \qquad $a(\alpha_1 + \alpha_2) = a \alpha_1 + a \alpha_2$;
\item[(ii)] $(a \cdot b) \cdot \alpha = a (b \alpha)$;
\item[(iii)] $1 \cdot \alpha = \alpha$.
\end{itemize}
\end{Def}

\begin{itemize}
\item 注: $f\big( (a, \alpha) \big) = a \cdot \alpha$.
\end{itemize}

\end{frame}

%------------------------------------------------

\section[例子]{例子}

%------------------------------------------------

\begin{frame}{Abel 群、$\mathbb{Z}$-模与线性空间}

\begin{eg}[Abel 群与 $\mathbb{Z}$-模]
$M$ 为一个 Abel 群, $R = \mathbb{Z}$
$\Longrightarrow$ $M$ 为一个 $\mathbb{Z}$-模; 反之亦然,
即 Abel 群 $\iff$ $\mathbb{Z}$-模.
\end{eg}

\pause

\begin{eg}[线性空间]
令 $R = F$ 为一个域, 一个 Abel 群 $V$ 为一个 $F$-模, 即 $V$ 为 $F$ 上的
一个线性空间 (有限维 or 无限维).
\end{eg}

\end{frame}

%------------------------------------------------

\begin{frame}{通过线性变换的 $F[X]$-模}

\begin{eg}[{$F[X]$}-模]
$R = F$ 为一个域, $V$ 为 $F$-模, $\sigma$ 为 $V \to V$ 的任一线性变换,
$F[X]$ 为一元多项式环, 此时 $V$ 也是 $F[X]$-模.
如果 $\forall\ \alpha \in V$, $f(X) \in F[X]$, 定义
\[
f(X) \cdot \alpha = f(\sigma) \cdot \alpha .
\]
\end{eg}

\end{frame}

%------------------------------------------------

\begin{frame}{理想作为模}

\begin{eg}[理想作为模]
子环 $N \subset R$ 为左理想 $\iff$ $N$ 为一个左 $R$-模;
\[
\text{右理想} \iff N \text{ 为一个 } R\text{-右模}.
\]
\end{eg}

\end{frame}

%------------------------------------------------

\begin{frame}{Poincar\'e 定理}

\begin{thm}[Poincar\'e 定理]
任何一个幺环都同构于某一个 Abel 群的自同态环.
\end{thm}

\startproof[bg=yellow, fg=red](证明)
令该 Abel 群 $M = (R, +)$.
\qed

\end{frame}

%------------------------------------------------

\begin{frame}{循环群与 $\mathbb{Z}_n^+$}

\begin{eg}[循环群与 $\mathbb{Z}_n^+$]
\begin{itemize}
\item[(1)] $M = \langle g \rangle$ 为一个无限循环群
  $\Longrightarrow$ $\operatorname{End}(M) = \mathbb{Z}$;
  $M = \langle g \rangle$ 为 $n$ 阶循环群
  $\Longrightarrow$ $\operatorname{End}(M) = \mathbb{Z}/\langle n \rangle$.
\item[(ii)] 令 $\mathbb{Z}_n^+ = \mathbb{Z}/\langle n \rangle$ (的加法群),
  \[
  \operatorname{Hom} \big( \mathbb{Z}^+, \mathbb{Z}_n^+ \big) = \mathbb{Z}_n^+ ;
  \qquad
  \operatorname{Hom} \big( \mathbb{Z}_n^+, \mathbb{Z}^+ \big) = 0 .
  \]
\end{itemize}
\end{eg}

\end{frame}

%------------------------------------------------

\begin{frame}{有理数加法群的自同态环}

\begin{eg}[有理数加法群]
令 $\mathbb{Q}^+$ 为有理数加法群, 则
$\operatorname{End}(\mathbb{Q}^+) \cong \mathbb{Q}$;
$\operatorname{Hom}(\mathbb{Q}^+, \mathbb{Z}^+) = 0$.
\end{eg}

\startproof[bg=yellow, fg=red]((iii) 的证明)
$\forall$ 给定的 $\sigma \in \operatorname{End}(\mathbb{Q}^+)$,
$\sigma(1) = x \in \mathbb{Q}$, 记为 $\sigma_x = \sigma$.
$\forall\ n \geqslant 1$, $\sigma_x \big( \tfrac{1}{n} \big) = \tfrac{x}{n}$,
因此, $\forall\ \tfrac{a}{b} \in \mathbb{Q}$, 有
$\sigma_x \big( \tfrac{a}{b} \big) = \tfrac{a}{b} x$
\[
\Longrightarrow\ \operatorname{End}(\mathbb{Q}^+) = \big\{ \sigma_x \;\big|\; x \in \mathbb{Q} \big\}
\]
定义 $x \mapsto \sigma_x$ 为环同构 $\mathbb{Q} \to \operatorname{End}(\mathbb{Q}^+)$.
\qed

\end{frame}

%------------------------------------------------

\section[模与同态]{模与环同态}

%------------------------------------------------

\begin{frame}{模与环同态}

$R$ 为一个幺环, $M$ 为 Abel 群, $\operatorname{End}(M)$ 为 $M$
的全自同态环.

\begin{prop}[模与环同态]
\begin{itemize}
\item $M$ 为 $R$-模 $\iff$ $\exists\ R \xrightarrow{\ \psi\ } \operatorname{End}(M)$
  为幺环同态;
\item $M$ 为 $R$-右模 $\iff$ $\exists\ R \xrightarrow{\ \psi\ } \operatorname{End}(M)$
  为幺环反同态.
\end{itemize}
\end{prop}

\end{frame}

%------------------------------------------------

\begin{frame}{模与环同态的证明}

\startproof[bg=yellow, fg=red](证明)
$M$ 为 $R$-模, 据定义, $\forall\ a \in R$, $\alpha \in M$, $a\alpha \in M$.
定义 $R \xrightarrow{\ \varphi\ } \operatorname{End}(M)$, by
$\forall\ a \in R$, $\varphi(a)(\alpha) = a \cdot \alpha$
\[
\Longrightarrow\ \varphi(a) \in \operatorname{End}(M),
\]
那么 $\varphi$ 为 $R \to \operatorname{End}(M)$ 的幺环同态
(都要由定义去验证).
\pause
反过来, $R \xrightarrow{\ \varphi\ } \operatorname{End}(M)$ 为幺环同态,
$\forall\ a \in R$, $\varphi(a) \in \operatorname{End}(M)$,
定义数乘 by, $\forall\ a \in R$, $\alpha \in M$,
$a \cdot \alpha = \varphi(a)(\alpha)$,
知 $M$ 为一个 $R$-模 (易验证).
\qed

\end{frame}

%------------------------------------------------

\begin{frame}{零化子}

若 $M$ 为 $R$ 模, $R \xrightarrow{\ \varphi_M\ } \operatorname{End}(M)$,
$\ker(\varphi_M)$ 为 $R$ 的一个理想.
\[
\ker(\varphi_M) = \big\{ a \in R \;\big|\; \varphi_M(a) = 0,\
\text{即 } \forall\ \alpha \in M,\ a \cdot \alpha = 0 \big\},
\]
称 $\ker(\varphi_M)$ 为 $M$ 作为 $R$-模的\alert{零化子}, 记为
$\operatorname{Ann}_R(M) = \ker(\varphi_M)$, 即
\[
\operatorname{Ann}_R(M)
= \big\{ a \in R \;\big|\; \forall\ \alpha \in M,\ a \cdot \alpha = 0 \big\} .
\]

\end{frame}

%------------------------------------------------

\section[结构传递]{模结构的传递与推论}

%------------------------------------------------

\begin{frame}{模结构的传递}

\begin{lem}[模结构的传递]
令 $R \xrightarrow{\ \sigma\ } R'$ 为幺环同态, 若 $M$ 为 $R'$-模,
那么 $M$ 必是 $R$-模;
反过来, 若 $R \xrightarrow{\ \sigma\ } R'$ 为满的环同态,
若 $M$ 为 $R$-模, 且 $\ker \sigma \subset \operatorname{Ann}_R(M)$,
则 $M$ 为 $R'$-模.
\end{lem}

\end{frame}

%------------------------------------------------

\begin{frame}{模结构传递的证明}

\startproof[bg=yellow, fg=red](证明)
$1^{\circ}$. $M$ 为 $R'$-模, 知
$R' \xrightarrow{\ \varphi\ } \operatorname{End}(M)$
$\Longrightarrow$ $\varphi \circ \sigma: R \to \operatorname{End}(M)$
为幺环同态, 那么即知 $M$ 为 $R$-模.
\pause
$2^{\circ}$. 性质: 若 $R \xrightarrow{\ \sigma\ } R'$ 为满的环同态,
则对于任一环同态 $R \xrightarrow{\ \tau\ } R''$,
且满足 $\ker \sigma \subset \ker \tau$, 则存在唯一的环同态
$R' \xrightarrow{\ \delta\ } R''$, 满足 $\tau = \delta \circ \sigma$
(习题).

\begin{center}
\begin{tikzpicture}[scale=0.9]
\node (R2) at (2.6, 1.8) {$R''$};
\node (R0) at (0, 0) {$R$};
\node (R1) at (5.2, 0) {$R'$};
\draw[->] (R0) -- node[above left] {$\tau$} (R2);
\draw[->] (R1) -- node[above right] {$\delta$} (R2);
\draw[->] (R0) -- node[below] {$\sigma$ (满)} (R1);
\node at (2.6, -1.0) {(满足 $\ker \sigma \subset \ker \tau$)};
\end{tikzpicture}
\end{center}
\qed

\end{frame}

%------------------------------------------------

\begin{frame}{推论}

\begin{cor}
若 $M$ 为一个 $R$-模, $N$ 为 $R$ 的任一理想,
且 $N \subset \operatorname{Ann}_R(M)$,
那么 $M$ 为一个 $R/N$-模.
\end{cor}

\end{frame}

%------------------------------------------------

\end{document}
